\chapter{Unit 3: Strings, Lists, Tuples and Dictionaries}

\section{Section A: Multiple Choice Practice Questions (60 Curated Problems)}

\noindent This section contains 60 rigorous multiple-choice questions extracted from university course examination assessments. Review each question, test edge-case syntax, and verify code execution mechanics.
\vspace{3mm}

\noindent\textbf{Q1.} \hfill {\small\color{LPUOrange}\textbf{[string a compound data type | Easy]}}\\[2pt]
Why is a string considered a compound data type in Python?

\begin{enumerate}[label=(\alph*),topsep=2pt,itemsep=1pt,partopsep=0pt,parsep=0pt]
    \item \relax Because it is made up of smaller pieces called characters
    \item \relax Because it always contains exactly one element that must be a letter of the alphabet from A to Z
    \item \relax Because it can only store numbers
    \item \relax Because it can be modified in place
\end{enumerate}
\vspace{2mm}

\noindent\textbf{Q2.} \hfill {\small\color{LPUOrange}\textbf{[length | Easy]}}\\[2pt]
Which function returns the number of characters in a string?

\begin{enumerate}[label=(\alph*),topsep=2pt,itemsep=1pt,partopsep=0pt,parsep=0pt]
    \item \relax \texttt{length()}
    \item \relax \texttt{len()}
    \item \relax \texttt{count()}
    \item \relax \texttt{size()}
\end{enumerate}
\vspace{2mm}

\noindent\textbf{Q3.} \hfill {\small\color{LPUOrange}\textbf{[accessing elements | Easy]}}\\[2pt]
What is the index of the first character in a Python string?

\begin{enumerate}[label=(\alph*),topsep=2pt,itemsep=1pt,partopsep=0pt,parsep=0pt]
    \item \relax \texttt{0}
    \item \relax \texttt{length}
    \item \relax \texttt{1}
    \item \relax \texttt{-1}
\end{enumerate}
\vspace{2mm}

\noindent\textbf{Q4.} \hfill {\small\color{LPUOrange}\textbf{[string slices | Easy]}}\\[2pt]
What is the result of \texttt{"python"[0:3]}?

\begin{enumerate}[label=(\alph*),topsep=2pt,itemsep=1pt,partopsep=0pt,parsep=0pt]
    \item \relax \texttt{"pyth"}
    \item \relax \texttt{"pyt"}
    \item \relax \texttt{"tho"}
    \item \relax \texttt{"python"}
\end{enumerate}
\vspace{2mm}

\noindent\textbf{Q5.} \hfill {\small\color{LPUOrange}\textbf{[string traversal | Easy]}}\\[2pt]
Which loop is commonly used to traverse each character of a string?

\begin{enumerate}[label=(\alph*),topsep=2pt,itemsep=1pt,partopsep=0pt,parsep=0pt]
    \item \relax \texttt{if} statement
    \item \relax \texttt{try} block
    \item \relax \texttt{for} loop
    \item \relax \texttt{with} statement
\end{enumerate}
\vspace{2mm}

\noindent\textbf{Q6.} \hfill {\small\color{LPUOrange}\textbf{[comparison | Easy]}}\\[2pt]
What does the expression \texttt{"apple" == "apple"} evaluate to?

\begin{enumerate}[label=(\alph*),topsep=2pt,itemsep=1pt,partopsep=0pt,parsep=0pt]
    \item \relax \texttt{False}
    \item \relax \texttt{True}
    \item \relax \texttt{None}
    \item \relax \texttt{0}
\end{enumerate}
\vspace{2mm}

\noindent\textbf{Q7.} \hfill {\small\color{LPUOrange}\textbf{[find function | Easy]}}\\[2pt]
What does \texttt{"hello".find("l")} return?

\begin{enumerate}[label=(\alph*),topsep=2pt,itemsep=1pt,partopsep=0pt,parsep=0pt]
    \item \relax \texttt{3}
    \item \relax \texttt{2}
    \item \relax \texttt{1}
    \item \relax \texttt{-1}
\end{enumerate}
\vspace{2mm}

\noindent\textbf{Q8.} \hfill {\small\color{LPUOrange}\textbf{[find function | Easy]}}\\[2pt]
What does the \texttt{find()} method return when the substring is not present?

\begin{enumerate}[label=(\alph*),topsep=2pt,itemsep=1pt,partopsep=0pt,parsep=0pt]
    \item \relax \texttt{None}
    \item \relax \texttt{-1}
    \item \relax An error is raised because the searched substring could not be located anywhere in the string
    \item \relax \texttt{0}
\end{enumerate}
\vspace{2mm}

\noindent\textbf{Q9.} \hfill {\small\color{LPUOrange}\textbf{[looping and counting | Easy]}}\\[2pt]
What is a common use of a counter variable while looping through a string?

\begin{enumerate}[label=(\alph*),topsep=2pt,itemsep=1pt,partopsep=0pt,parsep=0pt]
    \item \relax To convert the string to uppercase
    \item \relax To reverse the string
    \item \relax To count occurrences of a character
    \item \relax To delete characters from the string
\end{enumerate}
\vspace{2mm}

\noindent\textbf{Q10.} \hfill {\small\color{LPUOrange}\textbf{[list values | Easy]}}\\[2pt]
Which of the following correctly creates a list of three numbers?

\begin{enumerate}[label=(\alph*),topsep=2pt,itemsep=1pt,partopsep=0pt,parsep=0pt]
    \item \relax \texttt{(1, 2, 3)}
    \item \relax \texttt{[1, 2, 3]}
    \item \relax \texttt{\{1, 2, 3\}}
    \item \relax \texttt{<1, 2, 3>}
\end{enumerate}
\vspace{2mm}

\noindent\textbf{Q11.} \hfill {\small\color{LPUOrange}\textbf{[length | Easy]}}\\[2pt]
What does \texttt{len([10, 20, 30, 40])} return?

\begin{enumerate}[label=(\alph*),topsep=2pt,itemsep=1pt,partopsep=0pt,parsep=0pt]
    \item \relax \texttt{3}
    \item \relax \texttt{4}
    \item \relax \texttt{5}
    \item \relax \texttt{0}
\end{enumerate}
\vspace{2mm}

\noindent\textbf{Q12.} \hfill {\small\color{LPUOrange}\textbf{[membership | Easy]}}\\[2pt]
What does the expression \texttt{3 in [1, 2, 3]} evaluate to?

\begin{enumerate}[label=(\alph*),topsep=2pt,itemsep=1pt,partopsep=0pt,parsep=0pt]
    \item \relax \texttt{False}
    \item \relax \texttt{3}
    \item \relax \texttt{True}
    \item \relax \texttt{None}
\end{enumerate}
\vspace{2mm}

\noindent\textbf{Q13.} \hfill {\small\color{LPUOrange}\textbf{[operations | Easy]}}\\[2pt]
What is the result of \texttt{[1, 2] + [3, 4]}?

\begin{enumerate}[label=(\alph*),topsep=2pt,itemsep=1pt,partopsep=0pt,parsep=0pt]
    \item \relax \texttt{[1, 2, [3, 4]]}
    \item \relax \texttt{[4, 6]}
    \item \relax \texttt{[1, 2, 3, 4]}
    \item \relax \texttt{[3, 4, 1, 2]}
\end{enumerate}
\vspace{2mm}

\noindent\textbf{Q14.} \hfill {\small\color{LPUOrange}\textbf{[slices | Easy]}}\\[2pt]
What is the result of \texttt{[10, 20, 30, 40][1:3]}?

\begin{enumerate}[label=(\alph*),topsep=2pt,itemsep=1pt,partopsep=0pt,parsep=0pt]
    \item \relax \texttt{[20, 30, 40]}
    \item \relax \texttt{[20, 30]}
    \item \relax \texttt{[30, 40]}
    \item \relax \texttt{[10, 20]}
\end{enumerate}
\vspace{2mm}

\noindent\textbf{Q15.} \hfill {\small\color{LPUOrange}\textbf{[deletion | Easy]}}\\[2pt]
Which keyword is used to remove an item from a list by its index?

\begin{enumerate}[label=(\alph*),topsep=2pt,itemsep=1pt,partopsep=0pt,parsep=0pt]
    \item \relax \texttt{del}
    \item \relax \texttt{remove}
    \item \relax \texttt{erase}
    \item \relax \texttt{drop}
\end{enumerate}
\vspace{2mm}

\noindent\textbf{Q16.} \hfill {\small\color{LPUOrange}\textbf{[list and for loops | Easy]}}\\[2pt]
What is printed by \texttt{for x in [1, 2, 3]: print(x)}?

\begin{enumerate}[label=(\alph*),topsep=2pt,itemsep=1pt,partopsep=0pt,parsep=0pt]
    \item \relax The list \texttt{[1, 2, 3]}
    \item \relax The sum \texttt{6}
    \item \relax The length \texttt{3}
    \item \relax Each element \texttt{1}, \texttt{2}, \texttt{3} on separate lines
\end{enumerate}
\vspace{2mm}

\noindent\textbf{Q17.} \hfill {\small\color{LPUOrange}\textbf{[list parameters and nested list | Easy]}}\\[2pt]
What is a nested list?

\begin{enumerate}[label=(\alph*),topsep=2pt,itemsep=1pt,partopsep=0pt,parsep=0pt]
    \item \relax A list with only numbers
    \item \relax A list that contains one or more lists as elements
    \item \relax A list stored in a file
    \item \relax A list that cannot be changed
\end{enumerate}
\vspace{2mm}

\noindent\textbf{Q18.} \hfill {\small\color{LPUOrange}\textbf{[tuples | Easy]}}\\[2pt]
Which symbol pair is typically used to create a tuple?

\begin{enumerate}[label=(\alph*),topsep=2pt,itemsep=1pt,partopsep=0pt,parsep=0pt]
    \item \relax Parentheses \texttt{()}
    \item \relax Square brackets \texttt{[]}
    \item \relax Curly braces \texttt{\{\}}
    \item \relax Angle brackets \texttt{<>}
\end{enumerate}
\vspace{2mm}

\noindent\textbf{Q19.} \hfill {\small\color{LPUOrange}\textbf{[operations on tuples | Easy]}}\\[2pt]
What key property makes tuples different from lists?

\begin{enumerate}[label=(\alph*),topsep=2pt,itemsep=1pt,partopsep=0pt,parsep=0pt]
    \item \relax Tuples are immutable and cannot be changed after creation
    \item \relax Tuples use square brackets
    \item \relax Tuples can only hold strings
    \item \relax Tuples cannot be indexed
\end{enumerate}
\vspace{2mm}

\noindent\textbf{Q20.} \hfill {\small\color{LPUOrange}\textbf{[creating dictionary | Easy]}}\\[2pt]
Which of the following correctly creates a dictionary?

\begin{enumerate}[label=(\alph*),topsep=2pt,itemsep=1pt,partopsep=0pt,parsep=0pt]
    \item \relax \texttt{\{"a", 1, "b", 2\}}
    \item \relax \texttt{["a": 1, "b": 2]}
    \item \relax \texttt{\{"a": 1, "b": 2\}}
    \item \relax \texttt{("a", 1, "b", 2)}
\end{enumerate}
\vspace{2mm}

\noindent\textbf{Q21.} \hfill {\small\color{LPUOrange}\textbf{[string slices | Medium]}}\\[2pt]
What is the output of the following code? 

\begin{pycode}
s = "Programming"\\
print(s[2:8:2])
\end{pycode}

\begin{enumerate}[label=(\alph*),topsep=2pt,itemsep=1pt,partopsep=0pt,parsep=0pt]
    \item \relax \texttt{oai}
    \item \relax \texttt{rgamn}
    \item \relax \texttt{ogam}
    \item \relax \texttt{orm}
\end{enumerate}
\vspace{2mm}

\noindent\textbf{Q22.} \hfill {\small\color{LPUOrange}\textbf{[string traversal | Medium]}}\\[2pt]
What does the following loop print? 

\begin{pycode}
word = "abc"\\
for ch in word:\\
print(ch * 2, end="")
\end{pycode}

\begin{enumerate}[label=(\alph*),topsep=2pt,itemsep=1pt,partopsep=0pt,parsep=0pt]
    \item \relax \texttt{aaabbbccc}
    \item \relax \texttt{abcabc}
    \item \relax \texttt{aabbcc}
    \item \relax \texttt{abc}
\end{enumerate}
\vspace{2mm}

\noindent\textbf{Q23.} \hfill {\small\color{LPUOrange}\textbf{[find function | Medium]}}\\[2pt]
What is the result of the following code? 

\begin{pycode}
s = "banana"\\
print(s.find("na"))
\end{pycode}

\begin{enumerate}[label=(\alph*),topsep=2pt,itemsep=1pt,partopsep=0pt,parsep=0pt]
    \item \relax \texttt{4}
    \item \relax \texttt{-1}
    \item \relax \texttt{3}
    \item \relax \texttt{2}
\end{enumerate}
\vspace{2mm}

\noindent\textbf{Q24.} \hfill {\small\color{LPUOrange}\textbf{[comparison | Medium]}}\\[2pt]
What is the output of the following comparison? 

\begin{pycode}
print("apple" < "banana")
\end{pycode}

\begin{enumerate}[label=(\alph*),topsep=2pt,itemsep=1pt,partopsep=0pt,parsep=0pt]
    \item \relax \texttt{0}
    \item \relax \texttt{False}
    \item \relax \texttt{True}
    \item \relax \texttt{Error}
\end{enumerate}
\vspace{2mm}

\noindent\textbf{Q25.} \hfill {\small\color{LPUOrange}\textbf{[looping and counting | Medium]}}\\[2pt]
How many times does the letter \texttt{s} appear based on this code? 

\begin{pycode}
text = "mississippi"\\
count = 0\\
for c in text:\\
if c == "s":\\
count += 1\\
print(count)
\end{pycode}

\begin{enumerate}[label=(\alph*),topsep=2pt,itemsep=1pt,partopsep=0pt,parsep=0pt]
    \item \relax \texttt{3}
    \item \relax \texttt{2}
    \item \relax \texttt{4}
    \item \relax \texttt{5}
\end{enumerate}
\vspace{2mm}

\noindent\textbf{Q26.} \hfill {\small\color{LPUOrange}\textbf{[length | Medium]}}\\[2pt]
What is the output of the following code? 

\begin{pycode}
s = "Hello, World!"\\
print(len(s))
\end{pycode}

\begin{enumerate}[label=(\alph*),topsep=2pt,itemsep=1pt,partopsep=0pt,parsep=0pt]
    \item \relax \texttt{11}
    \item \relax \texttt{14}
    \item \relax \texttt{13}
    \item \relax \texttt{12}
\end{enumerate}
\vspace{2mm}

\noindent\textbf{Q27.} \hfill {\small\color{LPUOrange}\textbf{[string a compound data type | Medium]}}\\[2pt]
Which statement about Python strings is correct?

\begin{enumerate}[label=(\alph*),topsep=2pt,itemsep=1pt,partopsep=0pt,parsep=0pt]
    \item \relax Strings are stored as a list of integers
    \item \relax Strings can only store alphabetic characters
    \item \relax Strings are immutable, so \texttt{s[0] = "X"} raises an error
    \item \relax Strings are mutable and support item assignment
\end{enumerate}
\vspace{2mm}

\noindent\textbf{Q28.} \hfill {\small\color{LPUOrange}\textbf{[operations | Medium]}}\\[2pt]
What is the output of the following code? 

\begin{pycode}
result = [1, 2] + [3, 4] * 2\\
print(result)
\end{pycode}

\begin{enumerate}[label=(\alph*),topsep=2pt,itemsep=1pt,partopsep=0pt,parsep=0pt]
    \item \relax \texttt{[1, 2, 6, 8]}
    \item \relax \texttt{[1, 2, 3, 4, 1, 2, 3, 4]}
    \item \relax \texttt{[4, 8, 3, 4]}
    \item \relax \texttt{[1, 2, 3, 4, 3, 4]}
\end{enumerate}
\vspace{2mm}

\noindent\textbf{Q29.} \hfill {\small\color{LPUOrange}\textbf{[slices | Medium]}}\\[2pt]
What does the following code print? 

\begin{pycode}
nums = [10, 20, 30, 40, 50]\\
print(nums[-3:-1])
\end{pycode}

\begin{enumerate}[label=(\alph*),topsep=2pt,itemsep=1pt,partopsep=0pt,parsep=0pt]
    \item \relax \texttt{[20, 30]}
    \item \relax \texttt{[40, 50]}
    \item \relax \texttt{[30, 40, 50]}
    \item \relax \texttt{[30, 40]}
\end{enumerate}
\vspace{2mm}

\noindent\textbf{Q30.} \hfill {\small\color{LPUOrange}\textbf{[membership | Medium]}}\\[2pt]
What is the output of the following code? 

\begin{pycode}
data = [[1, 2], [3, 4]]\\
print(3 in data)
\end{pycode}

\begin{enumerate}[label=(\alph*),topsep=2pt,itemsep=1pt,partopsep=0pt,parsep=0pt]
    \item \relax \texttt{[3, 4]}
    \item \relax \texttt{True}
    \item \relax \texttt{Error}
    \item \relax \texttt{False}
\end{enumerate}
\vspace{2mm}

\noindent\textbf{Q31.} \hfill {\small\color{LPUOrange}\textbf{[deletion | Medium]}}\\[2pt]
What is the value of \texttt{nums} after this code runs? 

\begin{pycode}
nums = [1, 2, 3, 4, 5]\\
del nums[1:3]\\
print(nums)
\end{pycode}

\begin{enumerate}[label=(\alph*),topsep=2pt,itemsep=1pt,partopsep=0pt,parsep=0pt]
    \item \relax \texttt{[1, 3, 4, 5]}
    \item \relax \texttt{[1, 4, 5]}
    \item \relax \texttt{[2, 3, 4, 5]}
    \item \relax \texttt{[1, 2, 5]}
\end{enumerate}
\vspace{2mm}

\noindent\textbf{Q32.} \hfill {\small\color{LPUOrange}\textbf{[accessing elements | Medium]}}\\[2pt]
What does the following code print? 

\begin{pycode}
matrix = [[1, 2, 3], [4, 5, 6]]\\
print(matrix[1][2])
\end{pycode}

\begin{enumerate}[label=(\alph*),topsep=2pt,itemsep=1pt,partopsep=0pt,parsep=0pt]
    \item \relax \texttt{5}
    \item \relax \texttt{3}
    \item \relax \texttt{6}
    \item \relax \texttt{4}
\end{enumerate}
\vspace{2mm}

\noindent\textbf{Q33.} \hfill {\small\color{LPUOrange}\textbf{[list and for loops | Medium]}}\\[2pt]
What is the output of the following code? 

\begin{pycode}
nums = [1, 2, 3, 4]\\
total = 0\\
for n in nums:\\
total += n * n\\
print(total)
\end{pycode}

\begin{enumerate}[label=(\alph*),topsep=2pt,itemsep=1pt,partopsep=0pt,parsep=0pt]
    \item \relax \texttt{30}
    \item \relax \texttt{100}
    \item \relax \texttt{10}
    \item \relax \texttt{20}
\end{enumerate}
\vspace{2mm}

\noindent\textbf{Q34.} \hfill {\small\color{LPUOrange}\textbf{[list parameters and nested list | Medium]}}\\[2pt]
What does the following code print? 

\begin{pycode}
def modify(lst):\\
lst.append(99)\\
~\\
nums = [1, 2, 3]\\
modify(nums)\\
print(nums)
\end{pycode}

\begin{enumerate}[label=(\alph*),topsep=2pt,itemsep=1pt,partopsep=0pt,parsep=0pt]
    \item \relax \texttt{[99, 1, 2, 3]}
    \item \relax \texttt{[1, 2, 3, 99]}
    \item \relax \texttt{[1, 2, 3]}
    \item \relax \texttt{[99]}
\end{enumerate}
\vspace{2mm}

\noindent\textbf{Q35.} \hfill {\small\color{LPUOrange}\textbf{[tuples | Medium]}}\\[2pt]
What happens when the following code is executed? 

\begin{pycode}
t = (1, 2, 3)\\
t[0] = 10\\
print(t)
\end{pycode}

\begin{enumerate}[label=(\alph*),topsep=2pt,itemsep=1pt,partopsep=0pt,parsep=0pt]
    \item \relax Prints \texttt{(10, 2, 3)}
    \item \relax Prints \texttt{[10, 2, 3]}
    \item \relax Prints \texttt{(1, 2, 3)}
    \item \relax Raises a \texttt{TypeError}
\end{enumerate}
\vspace{2mm}

\noindent\textbf{Q36.} \hfill {\small\color{LPUOrange}\textbf{[operations on tuples | Medium]}}\\[2pt]
What is the output of the following code? 

\begin{pycode}
t = (1, 2) * 3\\
print(len(t))
\end{pycode}

\begin{enumerate}[label=(\alph*),topsep=2pt,itemsep=1pt,partopsep=0pt,parsep=0pt]
    \item \relax \texttt{5}
    \item \relax \texttt{6}
    \item \relax \texttt{2}
    \item \relax \texttt{3}
\end{enumerate}
\vspace{2mm}

\noindent\textbf{Q37.} \hfill {\small\color{LPUOrange}\textbf{[creating dictionary | Medium]}}\\[2pt]
What is the output of the following code? 

\begin{pycode}
d = dict([("a", 1), ("b", 2)])\\
print(d["b"])
\end{pycode}

\begin{enumerate}[label=(\alph*),topsep=2pt,itemsep=1pt,partopsep=0pt,parsep=0pt]
    \item \relax \texttt{1}
    \item \relax \texttt{2}
    \item \relax \texttt{b}
    \item \relax \texttt{("b", 2)}
\end{enumerate}
\vspace{2mm}

\noindent\textbf{Q38.} \hfill {\small\color{LPUOrange}\textbf{[adding-modifying-retrieving dictionary values | Medium]}}\\[2pt]
What does the following code print? 

\begin{pycode}
d = \{"x": 1\}\\
d["y"] = 2\\
d["x"] = 5\\
print(d)
\end{pycode}

\begin{enumerate}[label=(\alph*),topsep=2pt,itemsep=1pt,partopsep=0pt,parsep=0pt]
    \item \relax \texttt{\{'x': 1, 'y': 2\}}
    \item \relax \texttt{\{'x': 5\}}
    \item \relax \texttt{\{'x': 5, 'y': 2\}}
    \item \relax \texttt{\{'x': 1, 'y': 2, 'x': 5\}}
\end{enumerate}
\vspace{2mm}

\noindent\textbf{Q39.} \hfill {\small\color{LPUOrange}\textbf{[deleting items | Medium]}}\\[2pt]
What is the output of the following code? 

\begin{pycode}
d = \{"a": 1, "b": 2, "c": 3\}\\
del d["b"]\\
print(list(d.keys()))
\end{pycode}

\begin{enumerate}[label=(\alph*),topsep=2pt,itemsep=1pt,partopsep=0pt,parsep=0pt]
    \item \relax \texttt{['a', 'c']}
    \item \relax \texttt{['b']}
    \item \relax \texttt{['a', 'b', 'c']}
    \item \relax \texttt{['a', 'c', 'b']}
\end{enumerate}
\vspace{2mm}

\noindent\textbf{Q40.} \hfill {\small\color{LPUOrange}\textbf{[dictionary methods | Medium]}}\\[2pt]
What is the output of the following code? 

\begin{pycode}
d = \{"a": 1, "b": 2\}\\
print(d.get("c", 0))
\end{pycode}

\begin{enumerate}[label=(\alph*),topsep=2pt,itemsep=1pt,partopsep=0pt,parsep=0pt]
    \item \relax \texttt{Error}
    \item \relax \texttt{0}
    \item \relax \texttt{None}
    \item \relax \texttt{c}
\end{enumerate}
\vspace{2mm}

\noindent\textbf{Q41.} \hfill {\small\color{LPUOrange}\textbf{[string slices | Hard]}}\\[2pt]
What is the output of the following code? 

\begin{pycode}
s = "programming"\\
print(s[::-2])
\end{pycode}

\begin{enumerate}[label=(\alph*),topsep=2pt,itemsep=1pt,partopsep=0pt,parsep=0pt]
    \item \relax \texttt{gimrgm}
    \item \relax \texttt{gimargo}
    \item \relax \texttt{pormig}
    \item \relax \texttt{gnimargorp}
\end{enumerate}
\vspace{2mm}

\noindent\textbf{Q42.} \hfill {\small\color{LPUOrange}\textbf{[comparison | Hard]}}\\[2pt]
Which expression evaluates to \texttt{True} in Python 3?

\begin{enumerate}[label=(\alph*),topsep=2pt,itemsep=1pt,partopsep=0pt,parsep=0pt]
    \item \relax \texttt{"Z" < "A"}
    \item \relax \texttt{"apple" < "Apple"}
    \item \relax \texttt{"ABC" > "abc"}
    \item \relax \texttt{"Apple" < "apple"}
\end{enumerate}
\vspace{2mm}

\noindent\textbf{Q43.} \hfill {\small\color{LPUOrange}\textbf{[find function | Hard]}}\\[2pt]
What does the following print? 

\begin{pycode}
s = "mississippi"\\
print(s.find("ss", 3))
\end{pycode}

\begin{enumerate}[label=(\alph*),topsep=2pt,itemsep=1pt,partopsep=0pt,parsep=0pt]
    \item \relax \texttt{5}
    \item \relax \texttt{2}
    \item \relax \texttt{4}
    \item \relax \texttt{-1}
\end{enumerate}
\vspace{2mm}

\noindent\textbf{Q44.} \hfill {\small\color{LPUOrange}\textbf{[looping and counting | Hard]}}\\[2pt]
What is the value of \texttt{count} after this loop? 

\begin{pycode}
s = "aAbBcCaA"\\
count = 0\\
for ch in s:\\
if ch.lower() == 'a':\\
count += 1\\
print(count)
\end{pycode}

\begin{enumerate}[label=(\alph*),topsep=2pt,itemsep=1pt,partopsep=0pt,parsep=0pt]
    \item \relax \texttt{4}
    \item \relax \texttt{2}
    \item \relax \texttt{3}
    \item \relax \texttt{8}
\end{enumerate}
\vspace{2mm}

\noindent\textbf{Q45.} \hfill {\small\color{LPUOrange}\textbf{[string traversal | Hard]}}\\[2pt]
What is the output? 

\begin{pycode}
s = "abcde"\\
result = ""\\
for i in range(len(s)):\\
result += s[i] * i\\
print(result)
\end{pycode}

\begin{enumerate}[label=(\alph*),topsep=2pt,itemsep=1pt,partopsep=0pt,parsep=0pt]
    \item \relax \texttt{abbcccdddd}
    \item \relax \texttt{aabbbcccc}
    \item \relax \texttt{abcde}
    \item \relax \texttt{bccdddeeee}
\end{enumerate}
\vspace{2mm}

\noindent\textbf{Q46.} \hfill {\small\color{LPUOrange}\textbf{[operations | Hard]}}\\[2pt]
What is printed? 

\begin{pycode}
a = [1, 2, 3]\\
b = a * 2\\
b[0] = 99\\
print(a[0], b[0])
\end{pycode}

\begin{enumerate}[label=(\alph*),topsep=2pt,itemsep=1pt,partopsep=0pt,parsep=0pt]
    \item \relax \texttt{1 1}
    \item \relax \texttt{99 99}
    \item \relax \texttt{1 99}
    \item \relax \texttt{99 1}
\end{enumerate}
\vspace{2mm}

\noindent\textbf{Q47.} \hfill {\small\color{LPUOrange}\textbf{[list parameters and nested list | Hard]}}\\[2pt]
What is the output? 

\begin{pycode}
def modify(lst):\\
lst = lst + [4]\\
lst.append(5)\\
x = [1, 2, 3]\\
modify(x)\\
print(x)
\end{pycode}

\begin{enumerate}[label=(\alph*),topsep=2pt,itemsep=1pt,partopsep=0pt,parsep=0pt]
    \item \relax \texttt{[1, 2, 3, 5]}
    \item \relax \texttt{[1, 2, 3, 4]}
    \item \relax \texttt{[1, 2, 3]}
    \item \relax \texttt{[1, 2, 3, 4, 5]}
\end{enumerate}
\vspace{2mm}

\noindent\textbf{Q48.} \hfill {\small\color{LPUOrange}\textbf{[slices | Hard]}}\\[2pt]
What is the result? 

\begin{pycode}
lst = [0, 1, 2, 3, 4, 5]\\
lst[1:4] = [10, 20]\\
print(lst)
\end{pycode}

\begin{enumerate}[label=(\alph*),topsep=2pt,itemsep=1pt,partopsep=0pt,parsep=0pt]
    \item \relax \texttt{[0, 10, 20, 3, 4, 5]}
    \item \relax \texttt{[0, 10, 20, 4, 5]}
    \item \relax \texttt{[0, 10, 20, 2, 3, 4, 5]}
    \item \relax \texttt{[10, 20, 3, 4, 5]}
\end{enumerate}
\vspace{2mm}

\noindent\textbf{Q49.} \hfill {\small\color{LPUOrange}\textbf{[deletion | Hard]}}\\[2pt]
What is the output? 

\begin{pycode}
lst = [10, 20, 30, 40, 50]\\
del lst[1:4:2]\\
print(lst)
\end{pycode}

\begin{enumerate}[label=(\alph*),topsep=2pt,itemsep=1pt,partopsep=0pt,parsep=0pt]
    \item \relax \texttt{[10, 30, 40, 50]}
    \item \relax \texttt{[10, 30, 50]}
    \item \relax \texttt{[10, 40, 50]}
    \item \relax \texttt{[20, 40]}
\end{enumerate}
\vspace{2mm}

\noindent\textbf{Q50.} \hfill {\small\color{LPUOrange}\textbf{[list and for loops | Hard]}}\\[2pt]
What does this print? 

\begin{pycode}
matrix = [[1, 2], [3, 4], [5, 6]]\\
total = 0\\
for row in matrix:\\
total += row[-1]\\
print(total)
\end{pycode}

\begin{enumerate}[label=(\alph*),topsep=2pt,itemsep=1pt,partopsep=0pt,parsep=0pt]
    \item \relax \texttt{9}
    \item \relax \texttt{12}
    \item \relax \texttt{21}
    \item \relax \texttt{6}
\end{enumerate}
\vspace{2mm}

\noindent\textbf{Q51.} \hfill {\small\color{LPUOrange}\textbf{[membership | Hard]}}\\[2pt]
What is the output? 

\begin{pycode}
nested = [[1, 2], [3, 4]]\\
print(3 in nested, [3, 4] in nested)
\end{pycode}

\begin{enumerate}[label=(\alph*),topsep=2pt,itemsep=1pt,partopsep=0pt,parsep=0pt]
    \item \relax \texttt{False True}
    \item \relax \texttt{True False}
    \item \relax \texttt{False False}
    \item \relax \texttt{True True}
\end{enumerate}
\vspace{2mm}

\noindent\textbf{Q52.} \hfill {\small\color{LPUOrange}\textbf{[accessing elements | Hard]}}\\[2pt]
What is printed? 

\begin{pycode}
grid = [[0]*3]*3\\
grid[0][0] = 1\\
print(grid)
\end{pycode}

\begin{enumerate}[label=(\alph*),topsep=2pt,itemsep=1pt,partopsep=0pt,parsep=0pt]
    \item \relax \texttt{[[1, 0, 0], [0, 0, 0], [1, 0, 0]]}
    \item \relax \texttt{[[1, 0, 0], [1, 0, 0], [1, 0, 0]]}
    \item \relax \texttt{[[1, 1, 1], [0, 0, 0], [0, 0, 0]]}
    \item \relax \texttt{[[1, 0, 0], [0, 0, 0], [0, 0, 0]]}
\end{enumerate}
\vspace{2mm}

\noindent\textbf{Q53.} \hfill {\small\color{LPUOrange}\textbf{[tuples | Hard]}}\\[2pt]
Which of the following creates a tuple containing a single integer element?

\begin{enumerate}[label=(\alph*),topsep=2pt,itemsep=1pt,partopsep=0pt,parsep=0pt]
    \item \relax \texttt{t = tuple(5)}
    \item \relax \texttt{t = (5,)}
    \item \relax \texttt{t = (5)}
    \item \relax \texttt{t = 5}
\end{enumerate}
\vspace{2mm}

\noindent\textbf{Q54.} \hfill {\small\color{LPUOrange}\textbf{[operations on tuples | Hard]}}\\[2pt]
What is the output? 

\begin{pycode}
t = (1, [2, 3], 4)\\
t[1].append(5)\\
print(t)
\end{pycode}

\begin{enumerate}[label=(\alph*),topsep=2pt,itemsep=1pt,partopsep=0pt,parsep=0pt]
    \item \relax \texttt{(1, [2, 3, 5], 4)}
    \item \relax \texttt{(1, [2, 3], 4)}
    \item \relax \texttt{TypeError raised}
    \item \relax \texttt{(1, [2, 3], 4, 5)}
\end{enumerate}
\vspace{2mm}

\noindent\textbf{Q55.} \hfill {\small\color{LPUOrange}\textbf{[creating dictionary | Hard]}}\\[2pt]
What does the following produce? 

\begin{pycode}
d = dict(zip("abc", [1, 2, 3, 4]))\\
print(d)
\end{pycode}

\begin{enumerate}[label=(\alph*),topsep=2pt,itemsep=1pt,partopsep=0pt,parsep=0pt]
    \item \relax \texttt{\{'a': 1, 'b': 2, 'c': 3, None: 4\}}
    \item \relax \texttt{ValueError raised}
    \item \relax \texttt{\{'a': 1, 'b': 2, 'c': 3, 'd': 4\}}
    \item \relax \texttt{\{'a': 1, 'b': 2, 'c': 3\}}
\end{enumerate}
\vspace{2mm}

\noindent\textbf{Q56.} \hfill {\small\color{LPUOrange}\textbf{[adding-modifying-retrieving dictionary values | Hard]}}\\[2pt]
What is printed? 

\begin{pycode}
d = \{'x': 1\}\\
print(d.get('y', 0) + d.setdefault('y', 5))\\
print(d)
\end{pycode}

\begin{enumerate}[label=(\alph*),topsep=2pt,itemsep=1pt,partopsep=0pt,parsep=0pt]
    \item \relax \texttt{5} then \texttt{\{'x': 1\}}
    \item \relax \texttt{5} then \texttt{\{'x': 1, 'y': 5\}}
    \item \relax \texttt{0} then \texttt{\{'x': 1\}}
    \item \relax \texttt{10} then \texttt{\{'x': 1, 'y': 5\}}
\end{enumerate}
\vspace{2mm}

\noindent\textbf{Q57.} \hfill {\small\color{LPUOrange}\textbf{[deleting items | Hard]}}\\[2pt]
What is the output? 

\begin{pycode}
d = \{'a': 1, 'b': 2, 'c': 3\}\\
val = d.pop('b')\\
res = d.pop('z', -1)\\
print(val, res, len(d))
\end{pycode}

\begin{enumerate}[label=(\alph*),topsep=2pt,itemsep=1pt,partopsep=0pt,parsep=0pt]
    \item \relax \texttt{2 -1 2}
    \item \relax \texttt{2 None 2}
    \item \relax \texttt{2 -1 3}
    \item \relax \texttt{None -1 2}
\end{enumerate}
\vspace{2mm}

\noindent\textbf{Q58.} \hfill {\small\color{LPUOrange}\textbf{[dictionary methods | Hard]}}\\[2pt]
What does this print? 

\begin{pycode}
d = \{'a': 1, 'b': 2\}\\
d.update(\{'b': 20, 'c': 30\})\\
print(sum(d.values()))
\end{pycode}

\begin{enumerate}[label=(\alph*),topsep=2pt,itemsep=1pt,partopsep=0pt,parsep=0pt]
    \item \relax \texttt{51}
    \item \relax \texttt{53}
    \item \relax \texttt{33}
    \item \relax \texttt{23}
\end{enumerate}
\vspace{2mm}

\noindent\textbf{Q59.} \hfill {\small\color{LPUOrange}\textbf{[operations on dictionary | Hard]}}\\[2pt]
What is the output? 

\begin{pycode}
d = \{1: 'a', 2: 'b'\}\\
for k in list(d.keys()):\\
d[k*10] = d[k]\\
print(len(d))
\end{pycode}

\begin{enumerate}[label=(\alph*),topsep=2pt,itemsep=1pt,partopsep=0pt,parsep=0pt]
    \item \relax \texttt{4}
    \item \relax \texttt{6}
    \item \relax \texttt{RuntimeError raised}
    \item \relax \texttt{2}
\end{enumerate}
\vspace{2mm}

\noindent\textbf{Q60.} \hfill {\small\color{LPUOrange}\textbf{[length | Hard]}}\\[2pt]
What is printed? 

\begin{pycode}
s = "a\textbackslash\{\}tb\textbackslash\{\}nc"\\
print(len(s))
\end{pycode}

\begin{enumerate}[label=(\alph*),topsep=2pt,itemsep=1pt,partopsep=0pt,parsep=0pt]
    \item \relax \texttt{5}
    \item \relax \texttt{7}
    \item \relax \texttt{3}
    \item \relax \texttt{6}
\end{enumerate}
\vspace{2mm}


\newpage
\section{Section B: Comprehensive Subjective Problem with Detailed Analytical Rubric}

\begin{subjectiveq}{3}{1}
Design an industrial-grade string manipulation and nested data structure algorithm for inventory and text parsing:
\begin{enumerate}[label=(\alph*)]
    \item Write a function \texttt{word\_frequency\_matrix(corpus)} that cleans text (removing punctuation, case normalization), counts token frequency, and returns a sorted descending dictionary.
    \item Write a function \texttt{merge\_inventories(warehouse\_a, warehouse\_b)} that merges two nested dictionaries containing item codes, quantities, and unit prices, computing the aggregated inventory value and updating total quantities.
\end{enumerate}
\end{subjectiveq}

\begin{solbox}
\textbf{Python Solution:}
\begin{pycode}
import string\\
~\\
def word\_frequency\_matrix(corpus):\\
\quad clean\_tokens = corpus.lower().translate(str.maketrans("", "", string.punctuation)).split()\\
\quad freq = \{\}\\
\quad for tok in clean\_tokens:\\
\quad \quad freq[tok] = freq.get(tok, 0) + 1\\
\quad sorted\_freq = dict(sorted(freq.items(), key=lambda item: item[1], reverse=True))\\
\quad return sorted\_freq\\
~\\
def merge\_inventories(w1, w2):\\
\quad merged = \{\}\\
\quad all\_keys = set(w1.keys()) | set(w2.keys())\\
\quad for k in all\_keys:\\
\quad \quad q1 = w1.get(k, \{\}).get("qty", 0)\\
\quad \quad q2 = w2.get(k, \{\}).get("qty", 0)\\
\quad \quad p1 = w1.get(k, \{\}).get("price", 0.0)\\
\quad \quad p2 = w2.get(k, \{\}).get("price", 0.0)\\
\quad \quad price = max(p1, p2)\\
\quad \quad merged[k] = \{"qty": q1 + q2, "price": price, "total\_val": (q1 + q2) * price\}\\
\quad return merged
\end{pycode}
\textbf{Complexity Analysis:} Dictionary lookups and insertions operate in average $\mathcal{O}(1)$ time. Set union executes in $\mathcal{O}(|w_1| + |w_2|)$, ensuring complete $\mathcal{O}(N \log N)$ complexity dominated by dictionary sorting.
\end{solbox}

\begin{examinertip}
\begin{itemize}[leftmargin=1.5em, itemsep=2pt]
    \item \textbf{Algorithmic Formulation (3 Marks):} Correct problem decomposition, function signatures, and parameter documentation.
    \item \textbf{Syntactic Rigor \& Implementation (4 Marks):} Valid Python syntax, proper error handling, defensive input validation, and avoidance of anti-patterns.
    \item \textbf{Complexity \& Edge Analysis (3 Marks):} Accurate time and space asymptotic analysis $\mathcal{O}(n)$ and testing edge conditions (null inputs, zero division, empty collections).
\end{itemize}
\end{examinertip}
